% !TeX root = ../../tesis.tex

\section{Exportaciones \geojson{} Estáticas para Tableau}
\label{anx:tg-exports}

El script \texttt{utils/generate\_exports.py} genera los archivos \geojson{}
estáticos que Tableau carga como fuentes espaciales cuando el \textbf{GeoProxy}
no está activo —por ejemplo, al abrir el libro de trabajo en una máquina sin
acceso a la \textsc{api}. Su ejecución requiere que \textbf{VFTModel} esté
corriendo en \texttt{localhost:8000}.

El Cuadro~\ref{tab:tg-exports} describe los cinco archivos que produce el
script.

\begin{table}[htbp]
    \centering
    \caption[Exportaciones \geojson{} de Transport-gis para Tableau]{Archivos
    \geojson{} estáticos generados por \texttt{generate\_exports.py} y
    el indicador que representan.}
    \label{tab:tg-exports}
    \begin{tabularx}{\textwidth}{|l|l|X|}
        \hline
        \textbf{Archivo} & \textbf{Geometría} & \textbf{Indicador} \\
        \hline
        \texttt{df\_puntos.geojson}           & Point      & Factor de Desviación $DF$ — pares O-D muestreados \\
        \texttt{fc\_puntos.geojson}           & Point      & Fuerza Capilar $CF_h$ por nodo \\
        \texttt{b\_puntos.geojson}            & Point      & Centralidad de Intermediación $B(v)$ \\
        \texttt{cobertura\_estaciones.geojson} & Polygon    & Cobertura espacial $C$ a radio~800~m \\
        \texttt{df\_rutas\_lineas.geojson}     & LineString & Trayectorias O-D reales (Camino~B) \\
        \hline
    \end{tabularx}
    \fuentefigura{Elaboración propia con base en \texttt{generate\_exports.py}
    (Transport-gis-zmvm-mjg, 2026).}
\end{table}

El Camino~A (\textsc{rest}) genera los primeros cuatro archivos directamente
desde los endpoints de \textbf{VFTModel}:

\begin{lstlisting}[style=estiloBase, language=bash,
  caption={Ejecución de \texttt{generate\_exports.py} — Camino~A.},
  label=lst:tg-exports-cmda,
  basicstyle=\small\ttfamily]
cd Transport-gis-zmvm-mjg
python utils/generate_exports.py
\end{lstlisting}

Para el quinto archivo (\texttt{df\_rutas\_lineas.geojson}), que requiere
trayectorias completas como \texttt{LineString}, se utiliza el Camino~B:
exportar \texttt{routes\_raw.json} desde el cuaderno de VFTModel y pasarlo como
argumento:

\begin{lstlisting}[style=estiloBase, language=bash,
  caption={Ejecución con Camino~B: trayectorias lineales desde el
  cuaderno de VFTModel.},
  label=lst:tg-exports-cmdb,
  basicstyle=\small\ttfamily]
python utils/generate_exports.py --linestrings utils/data/routes_raw.json
\end{lstlisting}

El Código~\ref{lst:tg-exports-func} muestra la función \texttt{run\_camino\_a},
que orquesta la descarga de los cuatro indicadores vía \textsc{rest} en
secuencia.

\lstinputlisting[ style=estiloPython, firstline=41, lastline=68, caption={Función \texttt{run\_camino\_a}: descarga de cuatro indicadores desde \textbf{VFTModel} y exportación a \geojson{}.}, label=lst:tg-exports-func, breaklines=true, basicstyle=\tiny\ttfamily ]{Anexos/Transport-gis-mjg/utils/generate_exports.txt}
